% Generated draft: requires mathematical and source review.
\documentclass[11pt]{article}
\usepackage[margin=25mm]{geometry}
\usepackage{fontspec}
\setmainfont{Latin Modern Roman}
\usepackage{amsmath,amssymb,mathtools,unicode-math}
\setmathfont{Latin Modern Math}
\usepackage{enumitem,needspace,xurl,hyperref}
\setlength{\parindent}{0pt}
\setlength{\parskip}{5pt}
\setlength{\emergencystretch}{3em}
\newcommand{\N}{\mathbb{N}}
\newcommand{\Z}{\mathbb{Z}}
\newcommand{\Q}{\mathbb{Q}}
\newcommand{\R}{\mathbb{R}}
\newcommand{\C}{\mathbb{C}}
\newcommand{\F}{\mathbb{F}}
\newcommand{\A}{\mathbb{A}}
\newcommand{\PP}{\mathbb{P}}
\newcommand{\E}{\mathbb{E}}
\newcommand{\eps}{\varepsilon}
\newcommand{\dd}{\,\mathrm{d}}
\begin{document}
\section*{UnsolvedMath ID: 701042}
\textbf{Source problem code:} catalogue entry
\subsection{Definitions and mathematical statement}
Let \(p\) be a prime. A pro-\(p\) group is a topological group that is an inverse limit of finite \(p\)-groups. For a pro-\(p\) group \(G\), define its topological lower central series by
\[
\gamma_1(G)=G,\qquad
\gamma_{i+1}(G)=\overline{\langle [x,y]:x\in\gamma_i(G),\ y\in G\rangle}\quad(i\geq 1),
\]
where \([x,y]=x^{-1}y^{-1}xy\) and the bar denotes closure in \(G\). Say that \(G\) has finite width if there exists an integer \(d\geq 0\) such that
\[
\lvert\gamma_i(G)/\gamma_{i+1}(G)\rvert\leq p^d
\qquad\text{for every integer }i\geq 1.
\]
The source calls this having finite width \(p^d\); its inequality does not require \(p^d\) to be the least possible bound.

Under the conventional definition, \(G\) is hereditarily just infinite if \(G\) is infinite and
\[
\text{for every open subgroup }H\leq G\text{ and every closed normal subgroup }N\trianglelefteq H,
\qquad N\neq\{1\}\Longrightarrow [H:N]<\infty.
\]
\textbf{Classification problem, conditional on this interpretation.} Describe the pro-\(p\) groups \(G\) that are hereditarily just infinite and have finite width, with \(p\) ranging over primes.

\textbf{Unresolved source qualification.} The supplied definition says that each open subgroup has no closed normal subgroup of infinite index, without excluding the trivial subgroup. The trivial subgroup is closed and normal and has infinite index in every infinite open subgroup of an infinite pro-\(p\) group. Consequently, the supplied wording is incompatible with the conventional definition unless the nontriviality qualification is inserted. The displayed classification problem uses that conventional interpretation explicitly and requires source confirmation.
\subsection{Short English statement}
Describe the infinite pro-\(p\) groups whose lower central factors have uniformly bounded orders and whose open subgroups have only finite-index nontrivial closed normal subgroups. The source omits the word ``nontrivial'' in the latter condition; this interpretation therefore requires confirmation. The supplied registry labels the problem open.
\subsection{Sources}
\begin{itemize}
\item Local UnsolvedMath registry, permanent ID \(701042\), source problem code \texttt{catalogue entry}, title \emph{Kourovka Notebook Problem 14.57}. Supplied upstream revision: \texttt{d371ee917ac9647fbb06cd6fc529deec6f3dda82}.
\item \emph{The Kourovka Notebook (Unsolved Problems in Group Theory)}, Problem 14.57, attributed to C. R. Leedham-Green, proposed in the 14th issue (1999). The supplied background identifies the 21st issue edition (2026), \texttt{arXiv:1401.0300v48}; the supplied collection description names E. I. Khukhro and V. D. Mazurov as editors.
\item The registry records status \texttt{open} and an update timestamp of \texttt{2026-10-16T14:00:00Z}. These are supplied registry metadata, not an independently verified assessment of the literature. The cited source has not been independently consulted for this draft.
\end{itemize}
\end{document}
